\section*{Chapter \ref{ch:g11:redox} --- Oxidation and Reduction}

\begin{solution}{exo:g11:redox:1}
\omterm{def:g11:redox:oxidant}{Oxidant}: a species able to gain \omterm{def:g9:inside-the-atom:nucleus}{electrons}. \omterm{def:g11:redox:oxidant}{Reductant}: a species able to
give \omterm{def:g9:inside-the-atom:nucleus}{electrons} away. \omterm{def:g11:redox:oxidation}{Oxidation}: a loss of \omterm{def:g9:inside-the-atom:nucleus}{electrons}. \omterm{def:g11:redox:oxidation}{Reduction}: a gain
of \omterm{def:g9:inside-the-atom:nucleus}{electrons}.
\end{solution}

\begin{solution}{exo:g11:redox:2}
\ce{Fe^{2+}}/\ce{Fe}; \ce{Ag+}/\ce{Ag}; \ce{I2}/\ce{I-};
\ce{Fe^{3+}}/\ce{Fe^{2+}}.
\end{solution}

\begin{solution}{exo:g11:redox:3}
\ce{Ag+ + e- <=> Ag}; \ce{Al^{3+} + 3e- <=> Al};
\ce{Fe^{3+} + e- <=> Fe^{2+}}; \ce{Cl2 + 2e- <=> 2Cl-}.
\end{solution}

\begin{solution}{exo:g11:redox:4}
An \omterm{def:g11:redox:oxidation}{oxidation} (\omterm{def:g9:inside-the-atom:nucleus}{electrons} are lost); iron acts as a \omterm{def:g11:redox:oxidant}{reductant}.
\end{solution}

\begin{solution}{exo:g11:redox:5}
Zinc is oxidised and is the \omterm{def:g11:redox:oxidant}{reductant}; \ce{Cu^{2+}} is reduced and is
the \omterm{def:g11:redox:oxidant}{oxidant}. Two \omterm{def:g9:inside-the-atom:nucleus}{electrons} per zinc \omterm{def:g7:atoms-and-molecules:atom}{atom}.
\end{solution}

\begin{solution}{exo:g11:redox:6}
Start from \ce{H2O2 -> H2O}: % ce-unbalanced-ok (step 1 of the method)
oxygen: two O on the left, so \ce{2H2O} on the right; hydrogen: 2 on
the left, 4 on the right, so \ce{2H+} on the left; charge: $+2$ on the
left, 0 on the right, so two \omterm{def:g9:inside-the-atom:nucleus}{electrons} on the left:
\ce{H2O2 + 2H+ + 2e- <=> 2H2O}.
\end{solution}

\begin{solution}{exo:g11:redox:7}
\ce{I2 + 2e- -> 2I-} and \ce{2S2O3^{2-} -> S4O6^{2-} + 2e-}; two
\omterm{def:g9:inside-the-atom:nucleus}{electrons} each, so \ce{I2 + 2S2O3^{2-} -> 2I- + S4O6^{2-}}.
\end{solution}

\begin{solution}{exo:g11:redox:8}
\ce{Cu -> Cu^{2+} + 2e-} and $2\times$ (\ce{Ag+ + e- -> Ag}):
\ce{Cu + 2Ag+ -> Cu^{2+} + 2Ag}. The \ce{Cu^{2+}} \omterm{def:g9:ions:ion}{ions} formed are blue.
\end{solution}

\begin{solution}{exo:g11:redox:9}
\ce{C6H8O6 -> C6H6O6 + 2H+ + 2e-} and \ce{I2 + 2e- -> 2I-}:
\ce{C6H8O6 + I2 -> C6H6O6 + 2I- + 2H+}. Vitamin C is the \omterm{def:g11:redox:oxidant}{reductant}.
\end{solution}

\begin{solution}{exo:g11:redox:10}
\ce{Al -> Al^{3+} + 3e-} gives 3 \omterm{def:g9:inside-the-atom:nucleus}{electrons}, \ce{Cu^{2+} + 2e- -> Cu}
takes 2; the smallest common number is 6, so the first is multiplied by
2 and the second by 3: \ce{2Al + 3Cu^{2+} -> 2Al^{3+} + 3Cu}.
\end{solution}

\begin{solution}{exo:g11:redox:11}
\ce{H2O2 + 2H+ + 2e- -> 2H2O} and $2\times$ (\ce{Fe^{2+} -> Fe^{3+} + e-}):
\ce{H2O2 + 2H+ + 2Fe^{2+} -> 2H2O + 2Fe^{3+}}.
\end{solution}

\begin{solution}{exo:g11:redox:12}
$n(\ce{Zn}) = 1.30/65.4 = \qty{1.99e-2}{mol}$. One copper \omterm{def:g7:atoms-and-molecules:atom}{atom} is
deposited per zinc \omterm{def:g7:atoms-and-molecules:atom}{atom} oxidised, so $n(\ce{Cu}) = \qty{1.99e-2}{mol}$
and $m(\ce{Cu}) = 1.99\times 10^{-2} \times 63.5 = \qty{1.26}{g}$.
\end{solution}

\begin{solution}{exo:g11:redox:13}
\ce{Cr2O7^{2-} + 14H+ + 6e- -> 2Cr^{3+} + 7H2O} and
$6\times$ (\ce{Fe^{2+} -> Fe^{3+} + e-}):
\ce{Cr2O7^{2-} + 14H+ + 6Fe^{2+} -> 2Cr^{3+} + 6Fe^{3+} + 7H2O}.
Six \ce{Fe^{2+}} per dichromate \omterm{def:g9:ions:ion}{ion}: $6 \times 0.010 = \qty{0.060}{mol}$.
\end{solution}

\begin{solution}{exo:g11:redox:14}
\ce{2ClO- + 4H+ + 2e- -> Cl2 + 2H2O} and \ce{2Cl- -> Cl2 + 2e-}; adding
and dividing by 2: \ce{ClO- + Cl- + 2H+ -> Cl2 + H2O}. An acid supplies
the \ce{H+} \omterm{def:g9:ions:ion}{ions}, and the reaction releases dichlorine, a toxic gas:
hence the warning never to \omterm{def:g2:mixing-and-dissolving:mix}{mix} the two products.
\end{solution}

\begin{solution}{exo:g11:redox:15}
\omterm{def:g9:inside-the-atom:nucleus}{Electrons} cannot travel freely through the solution
(\cref{prop:g11:redox:no-free-electrons}): those lost by zinc \omterm{def:g7:atoms-and-molecules:atom}{atoms} can
only be taken by a \ce{Cu^{2+}} \omterm{def:g9:ions:ion}{ion} that touches the \omterm{def:g9:periodic-table-first-look:metal}{metal}. The copper
\omterm{def:g7:atoms-and-molecules:atom}{atom} is therefore formed at the surface of the strip, and stays there.
\end{solution}

\begin{solution}{pb:g11:redox:1}
\textbf{1.} \ce{Cr2O7^{2-}}/\ce{Cr^{3+}} and \ce{CH3COOH}/\ce{CH3CH2OH}.

\textbf{2.} Ethanol is oxidised: it is the \omterm{def:g11:redox:oxidant}{reductant}.

\textbf{3.} \ce{Cr2O7^{2-} + 14H+ + 6e- <=> 2Cr^{3+} + 7H2O}.

\textbf{4.} Carbon is balanced (two on each side). Oxygen: two on the
left, one on the right, so \ce{H2O} on the right. Hydrogen: four on the
left, $6 + 2 = 8$ on the right, so \ce{4H+} on the left. Charge: $+4$ on
the left, 0 on the right, so four \omterm{def:g9:inside-the-atom:nucleus}{electrons} on the left:
\ce{CH3COOH + 4H+ + 4e- <=> CH3CH2OH + H2O}.

\textbf{5.} Twelve \omterm{def:g9:inside-the-atom:nucleus}{electrons}: $2\times$ the first, $3\times$ the second
reversed:
\ce{2Cr2O7^{2-} + 3CH3CH2OH + 16H+ -> 4Cr^{3+} + 3CH3COOH + 11H2O}.

\textbf{6.} Orange \ce{Cr2O7^{2-}} \omterm{def:g9:ions:ion}{ions} are reduced to green
\ce{Cr^{3+}} \omterm{def:g9:ions:ion}{ions} wherever ethanol reaches them.

\textbf{7.} $M = 2\times 39.1 + 2\times 52.0 + 7\times 16.0 =
\qty{294.2}{g/mol}$.

\textbf{8.} $n = 2.0\times 10^{-3}/294.2 = \qty{6.8e-6}{mol}$.

\textbf{9.} Three ethanol \omterm{def:g7:atoms-and-molecules:molecule}{molecules} per two dichromate \omterm{def:g9:ions:ion}{ions}:
$n = \tfrac{3}{2}\times 6.8\times 10^{-6} = \qty{1.0e-5}{mol}$.

\textbf{10.} $M(\ce{C2H6O}) = \qty{46.0}{g/mol}$, so
$m = 1.02\times 10^{-5}\times 46.0 = \qty{4.7e-4}{g}$, about
\qty{0.47}{mg}.

\textbf{11.} $0.25/0.47 \approx 0.53$: about \qty{53}{\percent} of the
dichromate.

\textbf{12.} $0.53 \times 10 \approx \qty{5.3}{cm}$.

\textbf{13.} $\qty{0.25}{mg} \times 2100 = \qty{525}{mg}$ per litre of
blood, about \qty{0.53}{g/L}.

\textbf{14.} The breath meets the grains at the inlet first; the
dichromate there is in excess and removes all the ethanol, so the
reaction is complete in a first stretch of the tube, which turns fully
green, before any ethanol reaches further.

\textbf{15.} The reaction is total: the green \ce{Cr^{3+}} \omterm{def:g9:ions:ion}{ions} do not go
back to dichromate, so a used tube cannot measure again.

% ledger: ghs:K2Cr2O7 (H350 among its hazard statements)
\textbf{16.} It is acutely toxic (GHS06) and a serious health hazard
(GHS08: its hazard statements include ``may cause cancer''), as well as
corrosive and toxic to aquatic life; a device used by the public must
not contain it.

\textbf{17.} \ce{O2 + 4H+ + 4e- -> 2H2O} and
\ce{CH3CH2OH + H2O -> CH3COOH + 4H+ + 4e-}:
\ce{CH3CH2OH + O2 -> CH3COOH + H2O}.

\textbf{18.} Four \omterm{def:g9:inside-the-atom:nucleus}{electrons}.

\textbf{19.} $n = 0.25\times 10^{-3}/46.0 = \qty{5.4e-6}{mol}$ of
ethanol; \omterm{def:g9:inside-the-atom:nucleus}{electrons} $4n = \qty{2.2e-5}{mol}$, that is
$2.2\times 10^{-5}\times 6.02\times 10^{23} = \num{1.3e19}$ \omterm{def:g9:inside-the-atom:nucleus}{electrons}.

\textbf{20.} $Q = 1.3\times 10^{19}\times 1.60\times 10^{-19} \approx
\qty{2.1}{C}$. Each ethanol \omterm{def:g7:atoms-and-molecules:molecule}{molecule} sends exactly four \omterm{def:g9:inside-the-atom:nucleus}{electrons}
through the wire, so the charge is proportional to the amount of
ethanol: measuring it measures the ethanol.
\end{solution}
